%==============================================================================
%  LISTE DES ACRONYMES (titre à gauche, liste sans tableau)
%==============================================================================
\markboth{Liste des acronymes}{Liste des acronymes}
\thispagestyle{plain}

\vspace*{1.2cm}
\noindent{\Huge\bfseries Liste des acronymes}
\vspace{1.1em}

\noindent
AU : Action Unit, unité d'action faciale (FACS).\\[0.45em]
BERT : Bidirectional Encoder Representations from Transformers,
  encodeur bidirectionnel pré-entraîné pour le texte.\\[0.45em]
CNN : Convolutional Neural Network, réseau de neurones convolutif.\\[0.45em]
CoI : Community of Inquiry, communauté d'enquête.\\[0.45em]
\textsc{DAiSEE} : Dataset for Affective States in E-Environments,
  corpus public d'états affectifs en environnement d'apprentissage.\\[0.45em]
FACS : Facial Action Coding System, système de codage des actions faciales.\\[0.45em]
FER : Facial Expression Recognition, reconnaissance d'expressions faciales.\\[0.45em]
GPT : Generative Pre-trained Transformer, transformeur génératif pré-entraîné.\\[0.45em]
ICAP : Interactive, Constructive, Active, Passive.\\[0.45em]
\textsc{IEMOCAP} : Interactive Emotional Dyadic Motion Capture,
  corpus de capture de mouvement dyadique émotionnel.\\[0.45em]
LLM : Large Language Model, grand modèle de langue.\\[0.45em]
MOOC : Massive Open Online Course, cours en ligne ouvert et massif.\\[0.45em]
RAG : Retrieval-Augmented Generation, génération augmentée par récupération.\\[0.45em]
SER : Speech Emotion Recognition, reconnaissance d'émotions de la parole.\\[0.45em]
TALN : Natural Language Processing, traitement automatique du langage naturel.\\[0.45em]
ViT : Vision Transformer, transformeur de vision.\\[0.45em]
WavLM : Wave Language Model, encodeur vocal auto-supervisé.
